\section{Introduction}
\label{sec:introduction}

Large language models (LLMs)~\cite{brown2020language} have emerged as a transformative
technology in artificial intelligence (AI). Recent advancements in LLMs have
significantly improved their capability. LLMs have demonstrated tremendous
potential in a wide range of domains, such as machine translation, text
summarization, and conversational agents~\cite{chatgpt}. As a company serving billions
of users, we have been aggressively integrating AI into our products, and we are
putting LLMs as a high priority to shape the future of our products.

Training LLMs is a daunting task that requires enormous computation resources.
The scaling law~\cite{kaplan2020scaling} dictates that the model size and the
training data size are critical factors that determine the model capability. To
achieve state-of-the-art model capability, many efforts have been devoted to
train large models with hundreds of billions or even trillions of parameters on
hundreds of billions or even trillions of tokens. For example, GPT-3~\cite{floridi2020gpt}
has 175 billion parameters and PaLM~\cite{chowdhery2022palm} has 540 billion parameters. Major
players in this field build large-scale AI clusters with tens of thousands of
GPUs to train LLMs.


Scaling LLM training to tens of thousands of GPUs brings unprecedented
challenges. As AI has been at the core of many of our products, we have extensive
experience in training deep neural networks (DNNs). Yet, training a model like
ResNet~\cite{he2016deep} only takes tens or hundreds of GPUs.
Compared to these models,
the scale of training LLMs
is unparallel. While we are not new to building and operating
large-scale GPU clusters, these clusters are normally shared by many training
jobs. Now, in the context of LLM training, a single job is occupying tens of
thousands of GPUs and taking all the resources. The sheer scale of
LLM training introduces two specific challenges from a systems perspective.

The first challenge is to achieve high training efficiency at scale. Model FLOPs
utilization (MFU) is the ratio of the observed throughput to the theoretical
maximum throughput assuming 100\% of peak FLOPs~\cite{narayanan2021efficient}.
It is a standard metric to evaluate training efficiency that directly translates
to end-to-end training speed. LLM training is not embarrassingly parallel. To
train an LLM, the model is split across GPUs and the GPUs heavily communicate
with each other to make progress. Besides communication, other factors such as
operator optimization, data preprocessing and GPU memory consumption also
contribute significantly to MFU.

The second challenge is to achieve high training stability at scale, i.e.,
maintaining high training efficiency throughout the training process. Stability
is particularly important from a production perspective, as LLMs take a long
time to train. Training an LLM with one trillion tokens can take weeks. The
scale and time are orders of magnitude larger than those of regular DNN training
jobs. Failures and stragglers are the norm rather than the exception for LLM
training. At such a scale, the consequences of failures and stragglers are
devastating. Failures are very expensive, and it is critical to reduce the
recovery time, given the large scale. A straggler not only affects its own work,
but slows down the entire job involving tens of thousands of GPUs. 

In this paper, we present the design, implementation and engineering experience
of \sysname, a production system for training LLMs at scale. \sysname enables us
to scale LLM training to more than 10,000 GPUs. We are able to harness the power
of the massive number of GPUs to train LLMs with high training efficiency and
stability. In building and operating \sysname, we apply two systems principles:
algorithm-system co-design and \revision{in-depth observability}. 

\sysname is a specialized system tailored for LLM training. Algorithm-system
co-design is a key principle to maximize performance for specialized systems,
which has been applied widely in computer systems. We apply this
principle to \sysname in the context of LLM training with a full-stack approach
that spans all important system components. We make several modifications and
incorporate effective optimization techniques to the model architecture,
including parallel transformer block~\cite{chowdhery2022palm},
sliding window attention~\cite{beltagy2020longformer} and LAMB optimizer~\cite{You2020Large}. We leverage mixed parallelism strategies that combine data parallelism, pipeline parallelism, tensor
parallelism, and sequence parallelism. Importantly, we design custom techniques based
on the pattern of each parallelism strategy to maximize the overlapping between
communication and computation. We apply prefetching and
tree-based loading to optimize the data pipeline. We leverage non-blocking
asynchronous operations and eliminate global barriers for large-scale collective
communication group initialization. We design a custom network topology, reduce
ECMP hash conflicts, customize congestion control, and tune retransmit timeout
parameters for high network performance.

Stability problems including failures and stragglers in large-scale systems are
notoriously hard to diagnose and fix. Many hard stability issues only emerge at
large scale, which can stem from a wide range of software and hardware faults
deep in the stack. Manually identifying and resolving every single issue is infeasible
given the scale and complexity of the system. \revision{We apply the principle of in-depth observability to build a set of diagnosis tools. By 'in-depth observability', we mean a comprehensive monitoring and visualization strategy that penetrates beyond surface-level metrics to gather detailed, granular data across every component of the system stack, aiming to create a multidimensional view of system performance. The set of tools allows us to diagnose the system and identify root causes, by uncovering the intricate interactions and dependencies that contribute to stability issues.} We develop a robust training framework to automate fault
localization and recovery. We design heartbeat messages encapsulating various
forms of information to facilitate real-time anomaly detection and provide
early warnings. We implement a suite of diagnostic tests to identify nodes
causing disruptions. We optimize the checkpointing and recovery procedure to
reduce interruptions. To troubleshoot nuanced cases caused by stragglers, we develop a performance analysis tool to record fine-grained
CUDA events and generate system-wide heat-map and timeline trace from a distributed view, and develop a 3D parallel
training visualization tool to show data dependencies between ranks for diagnosis.


\sysname is deployed in our datacenters to train LLMs for our products. Over the
years, we have built several AI clusters with different size and hardware
configurations. Our largest AI cluster has over 10,000 GPUs. In terms of
training efficiency, \sysname achieves 55.2\% MFU when training a standard 175B
transformer model on 12,288 GPUs, providing an improvement of 1.34$\times$
compared to the state-of-the-art open-source training framework
Megatron-LM~\cite{shoeybi2020megatronlm}. In terms of model converge and
stability,
we show a real production run of \sysname that trains a proprietary
model with hundreds of billions of parameters on multi-trillion tokens for
several weeks. Over the weeks, the loss continues to converge, and \sysname
repairs and recovers the training process for over 100 times in presence of failures.
We also share our experience in diagnosing and fixing some intriguing problems. 
\revision{We are working on open-sourcing components that can benefit the community on GitHub\footnote{https://github.com/volcengine/veScale}.}
